% s-135-p135-m002; original PDF 135
\Aphorism{s-135-p135-m002}{191}
Es giebt Zustände gestörten Gleichgewichts, wo die Gegenwart nicht mehr Indifferenzpunkt zwischen Vergangenheit und Zukunft ist: eine gegenstandslose Unruhe, die Angst etwas zu versäumen oder zu verlieren, treibt vorwärts, die beklemmende Gewissheit schon begangener Fehlgriffe und Verirrungen schreckt zurück, hinter die letzten Erlebnisse, — man möchte zugleich flüchten und sich verkriechen, das Mass füllen und leeren, widerrufen und vollenden. Eine Nervenverstimmung ersten Ranges, eine Hyperästhesis und Lebensscheu, von der einzig eine eiskalte oder glühende Umarmung der Aussenwelt befreien kann: in solch einem Zustande müsste man in den Krieg ziehen oder ein Duell austragen oder heirathen.
\par

